\genetrace{} v0.1 is restricted to CS papers and to a single teacher
model (GPT-5.5). Cross-domain coverage (v0.2 biology and physics) and
multi-teacher consensus filtering are planned but not in scope for this
release. The six-field genome schema is inherited from our prior work
\ideaevolving{} and was not re-derived from first principles for this
paper; an ablation dropping individual fields (mechanism alone, niche
alone, etc.) would test whether all six are load-bearing. The five-mode
dynamics taxonomy is similarly inherited; alternative taxonomies (e.g.,
finer-grained transition types) are not explored.

\evoopd{} requires a teacher model with accessible per-token
log-probabilities. Open-weights teachers (Qwen3-14B-Thinking,
Llama-3-70B) satisfy this; closed APIs that return only top-$K$
log-probs (GPT-5.5 Black-Box, arXiv 2511.10643) work but with a known
truncation hit. The lineage-consistency signal requires the corpus to
provide multi-paper chains; tasks without natural lineage structure
(single-instance QA, open-domain reasoning) cannot benefit from
component (C) of the algorithm.

The contamination guard relies on the published \ideaevolving{}
\texttt{paper\_db} being a complete enumeration of \geneexam{} source
papers. If \geneexam{} draws from papers outside this list, our
denylist may be incomplete; the temporal cut (pre-2017) provides a
fallback guarantee. We measure Min-K\%++ on the test set as an
independent leakage check.

Reproducibility: all training runs are at LoRA scale and complete in
hours; replicating the full pipeline (SFT + \evoopd{}) requires 4$\times$
A100 80GB for $\sim$3 days. We provide the LoRA adapter weights, the
\genetrace{} v0.1 corpus, and the construction code under permissive
licenses to enable independent replication.
